% Options for packages loaded elsewhere
\PassOptionsToPackage{unicode}{hyperref}
\PassOptionsToPackage{hyphens}{url}
\documentclass[
  12pt,
]{ctexart}
\usepackage{xcolor}
\usepackage{amsmath,amssymb}
\setcounter{secnumdepth}{-\maxdimen} % remove section numbering
\usepackage{iftex}
\ifPDFTeX
  \usepackage[T1]{fontenc}
  \usepackage[utf8]{inputenc}
  \usepackage{textcomp} % provide euro and other symbols
\else % if luatex or xetex
  \usepackage{unicode-math} % this also loads fontspec
  \defaultfontfeatures{Scale=MatchLowercase}
  \defaultfontfeatures[\rmfamily]{Ligatures=TeX,Scale=1}
\fi
\usepackage{lmodern}
\ifPDFTeX\else
  % xetex/luatex font selection
\fi
% Use upquote if available, for straight quotes in verbatim environments
\IfFileExists{upquote.sty}{\usepackage{upquote}}{}
\IfFileExists{microtype.sty}{% use microtype if available
  \usepackage[]{microtype}
  \UseMicrotypeSet[protrusion]{basicmath} % disable protrusion for tt fonts
}{}
\makeatletter
\@ifundefined{KOMAClassName}{% if non-KOMA class
  \IfFileExists{parskip.sty}{%
    \usepackage{parskip}
  }{% else
    \setlength{\parindent}{0pt}
    \setlength{\parskip}{6pt plus 2pt minus 1pt}}
}{% if KOMA class
  \KOMAoptions{parskip=half}}
\makeatother
\usepackage{longtable,booktabs,array}
\usepackage{calc} % for calculating minipage widths
% Correct order of tables after \paragraph or \subparagraph
\usepackage{etoolbox}
\makeatletter
\patchcmd\longtable{\par}{\if@noskipsec\mbox{}\fi\par}{}{}
\makeatother
% Allow footnotes in longtable head/foot
\IfFileExists{footnotehyper.sty}{\usepackage{footnotehyper}}{\usepackage{footnote}}
\makesavenoteenv{longtable}
\setlength{\emergencystretch}{3em} % prevent overfull lines
\providecommand{\tightlist}{%
  \setlength{\itemsep}{0pt}\setlength{\parskip}{0pt}}
\usepackage{bookmark}
\IfFileExists{xurl.sty}{\usepackage{xurl}}{} % add URL line breaks if available
\urlstyle{same}
\hypersetup{
  hidelinks,
  pdfcreator={LaTeX via pandoc}}

\author{SCX}
\date{2026}

\begin{document}

\section{SCX 思想扩展 --- 基于 6 Agent
讨论的综合方案}\label{scx-ux601dux60f3ux6269ux5c55-ux57faux4e8e-6-agent-ux8ba8ux8bbaux7684ux7efcux5408ux65b9ux6848}

\begin{quote}
\textbf{TODO 部分已被 \texttt{SCX\_TODO\_2026-06-26.md} 替代}
(2026-06-26)

本文档保留仅用于所述 8 方向 + 3 阶段 roadmap 的策略参考。开发待办请查看
\texttt{SCX\_TODO\_2026-06-26.md}。
\end{quote}

\begin{quote}
生成：2026-06-26 \textbar{} 基于 356KB agent 输出的交叉分析
\end{quote}

\begin{center}\rule{0.5\linewidth}{0.5pt}\end{center}

\subsection{0. 现有基础}\label{ux73b0ux6709ux57faux7840}

\subsubsection{已完成的（v0.1.0）}\label{ux5df2ux5b8cux6210ux7684v0.1.0}

\begin{longtable}[]{@{}ll@{}}
\toprule\noalign{}
维度 & 状态 \\
\midrule\noalign{}
\endhead
\bottomrule\noalign{}
\endlastfoot
核心数学框架 (R\_m(s), SCX\_m(s), V(s)) & ✅ 完整定义 + 3 命题证明 \\
数据四分类 (Valuable/Redundant/Noisy/Expert-dep) & ✅ 规则引擎 \\
Python 包 (30 files, 6269 lines, 259 tests) & ✅ 全模块实现 \\
合成实验框架 & ✅ data\_generator + run\_experiment \\
数学根源追溯 & ✅ Kolmogorov → PAC-Bayes \\
竞争格局分析 & ✅ 6 维度覆盖，Expert Governance 是独有创新区 \\
SCX-Compress (加权 coreset) & ✅ 算法框架，缺理论证明 \\
\end{longtable}

\subsubsection{竞争格局中的独特位置}\label{ux7adeux4e89ux683cux5c40ux4e2dux7684ux72ecux7279ux4f4dux7f6e}

\begin{verbatim}
SCX 是唯一覆盖 6 个维度的方法：
  数据估值 + 数据选择 + 多专家 + 状态条件 + 噪声检测 + 可学习性

Expert Governance（蒸馏前专家规范化）是 SCX 的独有创新区
——没有已有工作在 MLIP 领域讨论跨专家冲突仲裁和域认证
\end{verbatim}

\begin{center}\rule{0.5\linewidth}{0.5pt}\end{center}

\subsection{1. 可发展的 8
个方向}\label{ux53efux53d1ux5c55ux7684-8-ux4e2aux65b9ux5411}

\subsubsection{方向 1：SCX-Compress 理论加固 🔴
P0}\label{ux65b9ux5411-1scx-compress-ux7406ux8bbaux52a0ux56fa-p0}

\textbf{现状}：算法已实现（weighted\_random/kcenter/herding），缺理论保证。

\textbf{发展内容}： - \textbf{Coreset 保真定理}：证明 SCX-Compress
的冗余分数 D(s) 与经验风险近似误差之间的 bound
\[\sup_{f \in \mathcal{F}} \left|\frac{1}{N}\sum_i \ell(f(x_i),y_i) - \frac{1}{N}\sum_j w_j \ell(f(x_j^*),y_j^*)\right| \leq \epsilon(D)\]
- \textbf{样本复杂度}：压缩到多少比例还能保持 95\% 精度？ - \textbf{与
Dataset Distillation 的等价性}：SCX-Compress
在什么条件下等价于梯度匹配蒸馏？

\textbf{实现}： - 在合成数据上验证 bound - CIFAR-100 上对比 coreset
baseline (K-Center, Herding, Random)

\begin{center}\rule{0.5\linewidth}{0.5pt}\end{center}

\subsubsection{方向 2：State Discovery 鲁棒性分析 🟡
P1}\label{ux65b9ux5411-2state-discovery-ux9c81ux68d2ux6027ux5206ux6790-p1}

\textbf{现状}：SCX 假设 state discovery 已完成，但真实场景中 state 是
latent 的。

\textbf{发展内容}： - \textbf{误指定分析}：如果状态划分有误差（merge
了两个不同状态、split 了一个同质状态），SCX 分类的结果会偏离多少？ -
\textbf{稳定性定理}：证明在什么条件下，状态划分的小扰动不会改变数据分类结果
- \textbf{自适应状态数}：不预设 K，根据 SCX 性能自动调整 -
\textbf{跨描述符稳定性}：用不同 φ（ACE/SOAP/MBTR）重复分析，量化 cluster
一致性

\textbf{实现}： - 合成实验中引入 controlled mis-specification - 添加
\texttt{StateDiscovery.stability\_analysis()} 方法

\begin{center}\rule{0.5\linewidth}{0.5pt}\end{center}

\subsubsection{方向 3：SCX 与 Influence Function 的融合 🟡
P1}\label{ux65b9ux5411-3scx-ux4e0e-influence-function-ux7684ux878dux5408-p1}

\textbf{现状}：竞争分析显示 Data Shapley 和 Influence Function
是主要对手，但 SCX 的 state-level 视角可以与它们互补。

\textbf{发展内容}： - \textbf{State-Conditioned Influence}：传统
influence 是 \(\mathcal{I}(x_i, x_j)\)，SCX 版本是
\(\mathcal{I}_s(x_i, x_j) = \mathcal{I}(x_i, x_j | x_i, x_j \in s)\) -
\textbf{两阶段估值}：SCX 粗选状态 → Influence/Shapley 细选样本 -
\textbf{与 LESS 的对比}：LESS 用梯度相似性选数据，SCX
用状态价值选数据------两者结合可以同时考虑''对模型影响大''和''状态覆盖不足''

\textbf{实现}： - 在 CIFAR-100 实验中加入 LESS baseline - 实现
\texttt{scx/valuation/influence.py}

\begin{center}\rule{0.5\linewidth}{0.5pt}\end{center}

\subsubsection{方向 4：Expert Governance Protocol 形式化 🔴
P0}\label{ux65b9ux5411-4expert-governance-protocol-ux5f62ux5f0fux5316-p0}

\textbf{现状}：竞争分析确认这是 SCX
的独有创新区------没有已有工作。但当前只是概念框架。

\textbf{发展内容}： - \textbf{Governance 协议}：形式化定义 expert
进入蒸馏前的必要步骤 1. Gauge Check → 2. Domain Certificate → 3.
Conflict Resolution → 4. Anchor Verification → 5. Distillation
Authorization - \textbf{Certification Bound}：给定 N 个 DFT
anchor，专家在状态 s 的可靠性估计误差不超过 ε 的概率 ≥ 1-δ -
\textbf{协议状态机}：每个 (expert, state) pair
经过协议后得到一个状态：\texttt{certified} /
\texttt{conditionally\_certified} / \texttt{uncertified} /
\texttt{rejected}

\textbf{实现}： - 在 AlN v3 数据上做 retrospective 验证（已有 DFT
标签可以作为 anchor） - 论文级输出：Protocol specification +
Certification theorem

\begin{center}\rule{0.5\linewidth}{0.5pt}\end{center}

\subsubsection{方向 5：SCX 通用 ML 实验 🟡
P1}\label{ux65b9ux5411-5scx-ux901aux7528-ml-ux5b9eux9a8c-p1}

\textbf{现状}：合成实验框架就绪，但缺真实数据实验。竞争分析中审稿人攻击点
4b 指出''合成实验不能证明在高维数据上有效''。

\textbf{发展内容}： - \textbf{CIFAR-100 + 5 ResNet-18}：每个 ResNet
在不同类别子集上训练，模拟 multi-expert 场景 - \textbf{UCI Tabular
Benchmark}：10+ 维真实数据，展示 SCX 不限于图像 - \textbf{对比}：Data
Shapley (via pyDVL), LOO, Random, Uncertainty, Diversity + SCX -
\textbf{关键指标}：相同标注预算下模型精度提升、数据分类准确率

\textbf{实现}： - 扩展现有 \texttt{experiments/ml\_benchmarks/} - 产出
2-3 张 Figure 级图表

\begin{center}\rule{0.5\linewidth}{0.5pt}\end{center}

\subsubsection{方向 6：SCX 自适应阈值 🟢
P2}\label{ux65b9ux5411-6scx-ux81eaux9002ux5e94ux9608ux503c-p2}

\textbf{现状}：四分类的阈值（error\_high=0.05, density\_high=0.05,
consistency\_high=0.7）是人为预设的。审稿人攻击点 4c
指出这是''参数化的数据描述''。

\textbf{发展内容}： -
\textbf{数据驱动的阈值选择}：在合成数据上扫描参数空间，找到使分类准确率最大的阈值
-
\textbf{自适应阈值}：根据数据分布自动调整------高维稀疏数据用更宽松的阈值
- \textbf{Bayesian 阈值}：将阈值建模为随机变量，用少量标注数据推断后验

\textbf{实现}： - 添加
\texttt{DataClassifier.auto\_calibrate(X,\ y\_labeled)} 方法 -
Sensitivity analysis: 阈值变化 ±20\% 对分类结果的影响

\begin{center}\rule{0.5\linewidth}{0.5pt}\end{center}

\subsubsection{方向 7：Online SCX 🔵
P3}\label{ux65b9ux5411-7online-scx-p3}

\textbf{现状}：SCX 是批处理模式------收集一批数据后统一分析。

\textbf{发展内容}： -
\textbf{增量状态更新}：新数据到来时，不重新聚类，而是更新现有状态的质心和半径
- \textbf{在线专家可靠性追踪}：用 exponential moving average 更新
R\_m(s) - \textbf{数据流版本的四分类}：实时判断 incoming sample 是
valuable/redundant/noisy/expert-dependent

\textbf{实现}： - \texttt{scx/core/online.py} --- OnlineSCXFramework -
合成数据流实验

\begin{center}\rule{0.5\linewidth}{0.5pt}\end{center}

\subsubsection{方向 8：SCX Benchmark Suite 🔵
P3}\label{ux65b9ux5411-8scx-benchmark-suite-p3}

\textbf{现状}：没有标准化的评估框架来比较 state-conditioned data
valuation 方法。

\textbf{发展内容}： - \textbf{Benchmark 设计}：5-10 个数据集（合成 + UCI
+ 图像 + MLIP），每个有预定义的 ground-truth state -
\textbf{评估指标}：state discovery accuracy, expert routing precision,
data classification F1, acquisition efficiency - \textbf{Baseline
集成}：Random, Uncertainty, Diversity, Data Shapley, LESS, Coreset

\textbf{实现}： - \texttt{experiments/benchmark/} 目录 - 标准化
API：\texttt{benchmark.run(method,\ dataset)\ →\ results}

\begin{center}\rule{0.5\linewidth}{0.5pt}\end{center}

\subsection{2. 优先级矩阵}\label{ux4f18ux5148ux7ea7ux77e9ux9635}

\begin{longtable}[]{@{}
  >{\raggedright\arraybackslash}p{(\linewidth - 10\tabcolsep) * \real{0.1200}}
  >{\raggedright\arraybackslash}p{(\linewidth - 10\tabcolsep) * \real{0.1600}}
  >{\raggedright\arraybackslash}p{(\linewidth - 10\tabcolsep) * \real{0.1800}}
  >{\raggedright\arraybackslash}p{(\linewidth - 10\tabcolsep) * \real{0.1800}}
  >{\raggedright\arraybackslash}p{(\linewidth - 10\tabcolsep) * \real{0.1800}}
  >{\raggedright\arraybackslash}p{(\linewidth - 10\tabcolsep) * \real{0.1800}}@{}}
\toprule\noalign{}
\begin{minipage}[b]{\linewidth}\raggedright
方向
\end{minipage} & \begin{minipage}[b]{\linewidth}\raggedright
优先级
\end{minipage} & \begin{minipage}[b]{\linewidth}\raggedright
理论深度
\end{minipage} & \begin{minipage}[b]{\linewidth}\raggedright
实验难度
\end{minipage} & \begin{minipage}[b]{\linewidth}\raggedright
论文贡献
\end{minipage} & \begin{minipage}[b]{\linewidth}\raggedright
时间估计
\end{minipage} \\
\midrule\noalign{}
\endhead
\bottomrule\noalign{}
\endlastfoot
1. SCX-Compress 理论 & 🔴 P0 & 高 & 低 & \textbf{Paper 4 核心定理} & 2-4
周 \\
4. Expert Governance Protocol & 🔴 P0 & 高 & 中 & \textbf{Paper 3
核心框架} & 4-6 周 \\
2. State Discovery 鲁棒性 & 🟡 P1 & 高 & 低 & Paper 4 鲁棒性章节 & 2-3
周 \\
3. SCX + Influence 融合 & 🟡 P1 & 中 & 中 & Paper 4 实验亮点 & 3-4 周 \\
5. 通用 ML 实验 & 🟡 P1 & 低 & 中 & \textbf{审稿人必需} & 3-5 周 \\
6. 自适应阈值 & 🟢 P2 & 低 & 低 & Paper 4 消融实验 & 1-2 周 \\
7. Online SCX & 🔵 P3 & 中 & 中 & 未来 Paper 5 & 4-8 周 \\
8. Benchmark Suite & 🔵 P3 & 低 & 高 & 社区影响力 & 6-10 周 \\
\end{longtable}

\begin{center}\rule{0.5\linewidth}{0.5pt}\end{center}

\subsection{3. Roadmap}\label{roadmap}

\begin{verbatim}
Phase 1 (1-2 月): 理论加固
  ├── SCX-Compress 定理 + 证明
  ├── State Discovery 鲁棒性定理
  └── Expert Governance Protocol 形式化

Phase 2 (2-3 月): 实验补强
  ├── CIFAR-100 + UCI 实验
  ├── SCX + Influence 对比实验
  └── 自适应阈值 sensitivity analysis

Phase 3 (3-6 月): 扩展
  ├── Online SCX 原型
  └── Benchmark Suite 设计
\end{verbatim}

\begin{center}\rule{0.5\linewidth}{0.5pt}\end{center}

\subsection{4.
与论文谱系的对应关系}\label{ux4e0eux8bbaux6587ux8c31ux7cfbux7684ux5bf9ux5e94ux5173ux7cfb}

\texttt{EGP\ Paper\ 1\ (egp/)} → \texttt{SCX-Theory} → \texttt{SCX-MLIP}
→ \texttt{SCX-Sim} → \texttt{SCX-Health}

\begin{longtable}[]{@{}
  >{\raggedright\arraybackslash}p{(\linewidth - 4\tabcolsep) * \real{0.2609}}
  >{\raggedright\arraybackslash}p{(\linewidth - 4\tabcolsep) * \real{0.3478}}
  >{\raggedright\arraybackslash}p{(\linewidth - 4\tabcolsep) * \real{0.3913}}@{}}
\toprule\noalign{}
\begin{minipage}[b]{\linewidth}\raggedright
方向
\end{minipage} & \begin{minipage}[b]{\linewidth}\raggedright
贡献给
\end{minipage} & \begin{minipage}[b]{\linewidth}\raggedright
具体位置
\end{minipage} \\
\midrule\noalign{}
\endhead
\bottomrule\noalign{}
\endlastfoot
1. SCX-Compress 理论 & \textbf{SCX-Theory (Paper 2)} & §4 Data Value →
§4.3 Compression Theory \\
2. State Discovery 鲁棒性 & \textbf{SCX-Theory (Paper 2)} & §3 State
Discovery → §3.4 Robustness Analysis \\
3. SCX + Influence 融合 & \textbf{SCX-MLIP (Paper 3)} & §6 Experiments →
§6.3 Comparison with Influence \\
4. Expert Governance & \textbf{SCX-MLIP (Paper 3)} & §3 Expert
Compilation Protocol \\
5. 通用 ML 实验 & \textbf{SCX-MLIP (Paper 3)} & §6.2 General ML
Experiments \\
6. 自适应阈值 & \textbf{SCX-MLIP (Paper 3)} & §5.1 Adaptive Threshold
Calibration \\
7. Online SCX & \textbf{SCX-Sim (Paper 4)} & 未来方向 \\
8. Benchmark Suite & 社区贡献 & 独立 release \\
\end{longtable}

\begin{center}\rule{0.5\linewidth}{0.5pt}\end{center}

\emph{本方案综合了 6 份 agent 输出（01-06，共
356KB），交叉分析竞争格局、数学基础、实现现状后提炼。}

\end{document}
